\section{\solaris Multiplayer Framework}
\label{sec:data_engine_details}

\subsection{Mineflayer Modifications}
To facilitate a wide range of Minecraft action recording, we modify the Mineflayer API to expose access to the most recently applied camera action in its physics module and extend its event system to send events on one-off semantic actions such as attacking, using, placing, and hotbar changes.

In addition, we introduce the following Mineflayer code to improve the quality of the collected data: 
\begin{itemize}[noitemsep, topsep=0pt]
  \item Make the bot correctly look at the face of the block when placing a new block.
  \item Add camera smoothing to all non-Pathfinder look commands.
\end{itemize}






\subsection{GPU Data Collection}

We initially experimented with \texttt{libOSMesa} for CPU-based headless rendering for the Minecraft clients (camera bots), however, videos in the open, \textit{Normal} world often appeared laggy, having up to four repeating frames at a 20 fps recording in graphics-intensive scenes, such as in a forest. Switching to GPU-based rendering produced consistently smooth videos without frame repetition. Additionally, switching to the GPU enabled the use of NVENC, Nvidia's hardware video encoder, which further reduced CPU workload as it had to encode screen captures via \texttt{ffmpeg}.
